\documentclass[12pt, letterpaper]{article}

\usepackage[utf8]{inputenc}
\title{Week 1}
\author{Raadhikka Gupta}
\date{\today}

\begin{document}
\maketitle

\section{Foundations in Generative AI}

\subsection{Generative AI process}
\begin{itemize}
    \item Data Input: AI model is given a large amount of data (images, sounds, text, and numbers)
    \item Training: AI model analyzes this data and looks for relationship between information. The \textbf{neural network} is trained on a dataset on the examples of \textbf{outputs}.
    \item Output: AI model generates new outputs by inputting a \textbf{random seed value}. It processes the seed value and generates new output based on the patterns learnt. 
\end{itemize}

\subsection{AI-assisted Development}
\begin{itemize}
    \item Building the Foundation
    \item Solving problems
    \item Safety Checks
    \item Documentation
    \item Refining and polishing
\end{itemize}
In all cases where AI is used, humans always double-check and ensure that the output produced by the AI is valid and doesn't cause any errors to the whole program.

\subsection{Large Language Models}
\begin{itemize}
    \item \textbf{Generative AI}: Training the model on TBs of data, (example) then asking it to generate the next word of a sentence/phrase
    \item \textbf{Foundational Model}: It is a type of Generative AI, that are trained on large amounts of data, and it is used as a base for a variety of different tasks
    \item \textbf{Large Language Model}: It is a type of Foundational Model, to understand and process human text. They are used for chatbots as they predict the next word in a sequence based on linguistic patterns. Examples: Claude, Gemini, GPT
    \item \textbf{Vision}: It is a type of Foundational Model, to analyze and understand visual data. It has pre-trained ML models for visual tasks like object classificaton, object detection, and image segmentation. 
    \item \textbf{Tuning}: The process of taking a trained model and training it more to make it more specific. Making it personalized for your use case, and giving it custom rules, tones, and data to its core behaviour.
    \item \textbf{Prompt Engineering}: To give instructions to an AI Model to generate an output.
\end{itemize}

\subsubsection{How LLMs are trained?}
\begin{itemize}
    \item \textbf{Exposed to dataset}: Given billions of pages from different sources that provides a foundational knowledge needed for training.
    \item \textbf{Pattern Recognition}: Identifies patterns in language like grammar, meaning, and context; by predicting the next word in a sequence of text.
    \item \textbf{Fine-tuning and iteration}: Repeated iterations and applications of algorithms allows it to generate output that are contextually appropriate and human-like.
\end{itemize}

\subsubsection{How do you work with a LLM?}
\begin{itemize}
    \item \textbf{Input Prompt}: A question, instruction, or statement.
    \item \textbf{Tokenisation}: Prompt is broken into \textbf{tokens}; words or subwords; to interpret and process the input carefully.
    \item \textbf{Contextual understanding}: Uses transformer architecture to analyze the prompt in context, to examine how words relate to each other.
    \item \textbf{Prediction}: Predicts the next most likely token based on the prompt, training data, and statistical patterns learnt.
    \item \textbf{Output generation}: Continues to predict tokens until reaches a stopping point, the output.
\end{itemize}

\subsection{Prompt}
\begin{itemize}
    \item A \textbf{prompt} is an input given to the Generative AI system to help it create a specific output.
\end{itemize}

\subsubsection{How to make a prompt work?}
\begin{itemize}
    \item \textbf{Clarity}: provides a clear task, good context, and avoids vague instructions.
    \item \textbf{Relevance}: matches the goal to complete, speaks to the correct audience, and uses a format that fits the task.
\end{itemize}

\subsubsection{How to write a prompt?}
\begin{itemize}
    \item \textbf{Role}: How the AI should act as and expertise it must use 
    \item \textbf{Context}: Background information 
    \item \textbf{Task}: What you want the AI to do
    \item \textbf{Output}: What the output should look like (format or structure)
    \item \textbf{Constraints}: Limits, such as length, tone, or required content
\end{itemize}

\end{document}